\title{Controlling Text-to-Image Diffusion by Orthogonal Finetuning}

\begin{abstract}
Large-scale text-to-image diffusion models demonstrate remarkable proficiency in producing photorealistic images from textual descriptions. However, devising effective strategies to guide or regulate these powerful models for a variety of downstream applications remains a significant unsolved problem. In response, we propose a systematic finetuning approach—Orthogonal Finetuning (OFT)—to tailor text-to-image diffusion models for specific downstream tasks. Distinct from prior approaches, OFT is theoretically guaranteed to maintain the hyperspherical energy, a property reflecting pairwise neuron relationships on the unit hypersphere. Our investigation shows that preserving this property is essential for retaining the semantic fidelity of image generation in text-to-image diffusion models. To further enhance the robustness of finetuning, we introduce Constrained Orthogonal Finetuning (COFT), which incorporates an additional constraint on the hypersphere's radius. Our study focuses on two key text-to-image finetuning tasks: (1) subject-driven generation, which aims to synthesize images depicting a particular subject—given a handful of images and a text prompt—in varied scenarios; and (2) controllable generation, where the model is equipped to integrate auxiliary control signals. Through empirical evaluation, we demonstrate that the OFT framework not only surpasses existing methods in generation quality but also accelerates convergence during training.
\end{abstract}

\section{Introduction}

Advancements in text-to-image diffusion models~\cite{saharia2022photorealistic,ramesh2022hierarchical,rombach2022high} have led to unprecedented results in generating high-fidelity images with text guidance. Nevertheless, textual conditioning alone often lacks the precision and granularity required for more exacting control over the resulting images. In this work, we address two specific scenarios in text-to-image generation:

\begin{itemize}[leftmargin=*,nosep]

    \item \textbf{Subject-driven generation}~\cite{ruiz2023dreambooth}: Here, the objective is to generate new images of a given subject—provided only a few example images—placed within different contexts, utilizing a text prompt as guidance.
    \item \textbf{Controllable generation}~\cite{zhang2023adding,mou2023t2i}: In this setting, the goal is to generate images conditioned not only on a text prompt but also on supplemental control inputs (such as canny edges or segmentation maps), ensuring the output conforms to these signals.
\end{itemize}

At their core, both tasks revolve around the challenge of effectively finetuning text-to-image diffusion models while maintaining the generative capabilities established during pretraining. The key criteria for a robust finetuning approach can be outlined as follows: (1) \emph{training efficiency}, which involves reducing the number of trainable parameters and training epochs, and (2) \emph{generalizability preservation}, which focuses on sustaining high-fidelity and diverse outputs. Common finetuning approaches include incrementally modifying neuron weights with a modest learning rate (\emph{e.g.}, \cite{ruiz2023dreambooth}), or incorporating a lightweight module with separately parameterized weights (\emph{e.g.}, \cite{hulora2022,zhang2023adding}). Although these methods are straightforward, they do not provide guarantees regarding the retention of pretraining generative quality. Moreover, there remains a notable absence of theoretical guidance for constructing optimal finetuning procedures or for determining appropriate hyperparameters, such as the duration of training. A central challenge is the unavailability of a reliable metric to assess how well the finetuned model preserves its original generative abilities. Present finetuning strategies often assume that minimizing the Euclidean distance between finetuned and pretrained parameters correlates with improved retention of pretraining strengths. This leads to the prevalent use of very small learning rates during adaptation. Nonetheless, while this assumption is sometimes valid, we contend that Euclidean proximity alone does not adequately reflect semantic preservation. Consequently, introducing a more structurally-informed measure for evaluating differences between the finetuned and pretrained models could substantially enhance both the maintenance of original generative performance and the stability of the finetuning process.

Building on previous findings that hyperspherical similarity effectively represents semantic relationships~\cite{liu2017deep,liu2018decoupled,chen2020angular}, we employ the notion of hyperspherical energy~\cite{liu2018learning} to model the pairwise interactions among neurons. Specifically, hyperspherical energy quantifies the aggregate hyperspherical similarities (\emph{e.g.}, cosine similarities) between all pairs of neurons within a layer, thus reflecting how uniformly the neurons are distributed on the unit hypersphere~\cite{liu2021learning}. We put forward the hypothesis that optimal finetuning should minimally alter the hyperspherical energy relative to the pretrained counterpart. While a straightforward solution is to incorporate a regularization term maintaining constant hyperspherical energy during finetuning, this approach lacks assurances that the energy gap will indeed be small. Consequently, we utilize the invariance property inherent to hyperspherical energy: applying a uniform orthogonal transformation across all neurons leaves the pairwise hyperspherical similarity unchanged. Based on this insight, we introduce Orthogonal Finetuning~(OFT), which adapts large-scale text-to-image diffusion models for downstream applications while preserving hyperspherical energy. The principal concept involves training a single, shared orthogonal transformation at the layer level, thereby maintaining the pairwise angular relationships among neurons. From another perspective, OFT can be interpreted as recalibrating the canonical coordinate system for the set of neurons within a layer. Additionally, considering that lower Euclidean distance between the parameters of the finetuned and pretrained models enhances the retention of pretraining performance, we propose a variant, Constrained Orthogonal Finetuning~(COFT), which confines the finetuned parameters within a hypersphere of fixed radius centered at the original neurons.

The rationale behind utilizing orthogonal transformation for finetuning neurons is partly inspired by the properties of the 2D Fourier transform, which enables the decomposition of an image into its magnitude and phase spectra. The phase spectrum, encoding the angular relationships between inputs and basis functions, is primarily responsible for retaining semantic content. For instance, one can reconstruct an image that maintains its semantics by combining its original phase spectrum with an arbitrary magnitude spectrum. This observation implies that modifying neuron orientations is central to achieving semantic adjustments in generated images—a critical objective in both subject-driven and controllable image generation. Nevertheless, allowing unrestricted changes in neuron orientation risks compromising the pretrained model’s generative abilities. To limit this freedom, we advocate for maintaining the inter-neuron angles, motivated by the premise that these angles play a pivotal role in encapsulating the knowledge within neural networks. Following this line of reasoning, we propose a layer-wise, shared orthogonal transformation among neurons in each layer, thereby ensuring the preservation of hyperspherical energy.

Additionally, our approach takes cues from orthogonal over-parameterized training~\cite{liu2021orthogonal}, which builds classification neural networks from the ground up by applying orthogonal transformations to networks initialized with random weights. The reason for this stems from the theoretical result that a randomly initialized neural network is expected to exhibit low hyperspherical energy, and the central objective of \cite{liu2021orthogonal} is to maintain this low energy throughout training—a property linked to improved generalization in classification tasks~\cite{liu2018learning,lin2020regularizing}. The findings in \cite{liu2021orthogonal} demonstrate the adequacy of orthogonal transformations for training classification models with strong generalization capabilities. However, our work is distinguished by its focus on finetuning text-to-image diffusion models to enhance controllability and downstream generative quality. OFT and \cite{liu2021orthogonal} diverge in two main respects. Firstly, unlike \cite{liu2021orthogonal}, which seeks to minimize hyperspherical energy, OFT strives to retain the pretrained model’s hyperspherical energy, safeguarding its underlying structural integrity during finetuning. Specifically, further reduction of hyperspherical energy during diffusion model finetuning could potentially disrupt intrinsic semantic structures. Secondly, OFT introduces an additional constraint by seeking to minimize deviations from the pretrained parameters, resulting in a constrained variant, whereas \cite{liu2021orthogonal} imposes no such limitation. Finetuning fundamentally involves balancing adaptability with stability, and we contend that our OFT framework successfully negotiates this trade-off. Our key contributions are summarized as follows:

\begin{itemize}[leftmargin=*,nosep]

    \item We introduce a new finetuning strategy—Orthogonal Finetuning—designed to enhance the controllability of text-to-image diffusion models. To further ensure finetuning stability, we present a constrained version that restricts the angular displacement from the original pretrained parameters.
    \item Unlike conventional finetuning approaches, OFT is constructed to preserve the model's hyperspherical energy during the adaptation process, a property that we empirically show to correlate with the retention of the pretrained model’s generative semantics.
    \item We evaluate OFT on both subject-driven and controllable image generation tasks. Through an extensive set of experiments, we show that OFT achieves notable gains compared to existing techniques, yielding superior generation fidelity, faster convergence, and enhanced stability during finetuning. Additionally, OFT is highly sample-efficient, demonstrating robust performance even when trained with considerably fewer images and epochs.
    \item For the controllable generation scenario, we develop a novel control consistency metric to quantitatively assess control fidelity. This metric works by extracting the control signal from the synthesized output and directly comparing it to the original control input, thereby further substantiating OFT’s improved controllability.
\end{itemize}

\section{Related Work}

\textbf{Text-to-image diffusion models}. Remarkable advancements~\cite{saharia2022photorealistic,ramesh2022hierarchical,rombach2022high,gu2022vector,nichol2022glide} in generating images from text descriptions have resulted largely from innovations in diffusion-based generative modeling~\cite{ho2020denoising,song2021score,song2021denoising,dhariwal2021diffusion} and in learning joint vision-language representations~\cite{su2020vlbert,lu2019vilbert,radford2021learning,wang2021simvlm,li2022blip,li2023blip,alayrac2022flamingo,schuhmann2022laion}. GLIDE~\cite{nichol2022glide} and Imagen~\cite{saharia2022photorealistic} operate by training diffusion models directly in pixel space, with GLIDE performing joint training of a text encoder and diffusion prior on aligned text-image pairs, whereas Imagen leverages a fixed, pretrained text encoder. Models such as Stable Diffusion~\cite{rombach2022high} and DALL-E2~\cite{ramesh2022hierarchical} apply diffusion to latent representations: Stable Diffusion utilizes VQ-GAN~\cite{esser2021taming} to derive its latent visual space, while DALL-E2 makes use of CLIP~\cite{radford2021learning} to build a shared embedding space unifying visual and textual inputs. Besides diffusion models, other generative paradigms, such as generative adversarial networks~\cite{reed2016generative,zhang2017stackgan,xu2018attngan,li2019controllable} and autoregressive models~\cite{ramesh2021zero,ding2021cogview,wu2022nuwa,yuscaling}, have been explored for this task. OFT is designed as a broadly applicable, model-agnostic finetuning method suitable for use with any text-to-image diffusion model.

\textbf{Subject-driven generation}. To prevent changes to the depicted subject, approaches in~\cite{avrahami2022blended,nichol2022glide} supplement the model input with user-supplied masks. Inversion-based techniques~\cite{choi2021ilvr,dhariwal2021diffusion,ramesh2022hierarchical,gal2022image} enable contextual modifications while maintaining the core subject. Meanwhile,~\cite{hertz2022prompt} facilitates local and global edits without the necessity for explicit masking. Nevertheless, these methods often struggle to maintain subject-specific features, particularly those related to identity. Pivotal Tuning~\cite{roich2022pivotal} finetunes the generator in the vicinity of an inverted latent representation, employing extra regularization to preserve identity. Similarly,~\cite{nitzan2022mystyle} constructs a personalized generative prior using a person’s facial image set. In~\cite{casanova2021instance}, multiple instance variants are produced, but at the expense of distinct instance characteristics. DreamBooth~\cite{ruiz2023dreambooth} updates the text-to-image diffusion model by introducing a unique token and a limited set of subject photos, utilizing reconstruction and class-preserving loss functions. Building upon the DreamBooth paradigm, OFT replaces straightforward low-rate finetuning with updates constrained to orthogonal transformations

\textbf{Controllable generation}. The image-to-image translation problem represents a special instance of controllable generation. Earlier approaches typically leverage conditional generative adversarial networks~\cite{isola2017image,zhu2017toward,wang2018high,park2019semantic,choi2018stargan}. Additionally, diffusion models have been explored for image-to-image translation tasks~\cite{saharia2022palette,voynov2022sketch,wang2022pretraining}. Recently, ControlNet~\cite{zhang2023adding} introduced a strategy of regulating a pretrained diffusion model by finetuning it with additional control cues, thus attaining state-of-the-art controllability in image generation. Likewise, the contemporaneous work T2I-Adapter~\cite{mou2023t2i} enhances controllability over generated images by similarly adapting a pretrained diffusion model through finetuning. Consistent with the setup in \cite{zhang2023adding,mou2023t2i}, our approach applies OFT for finetuning pretrained diffusion models, resulting in markedly improved controllability while requiring fewer data and reduced numbers of finetuning parameters. Importantly, OFT imposes no additional computational demands during inference.

\textbf{Model finetuning}. The practice of finetuning extensive pretrained models on downstream applications has seen increasing adoption in recent years~\cite{devlin2018bert,brown2020language,he2022masked}. In the context of finetuning, various adaptation methods (\emph{e.g.}, \cite{hulora2022,houlsby2019parameter,pfeiffer2020adapterfusion}) have been well explored within natural language processing. Among these, LoRA~\cite{hulora2022} is most closely related to OFT, as it posits a low-rank configuration for additive parameter modifications during adaptation. Distinctly, OFT employs a multiplicative update via a layer-wise shared orthogonal transformation, thereby modulating neuron weights. This method is mathematically guaranteed to maintain the pairwise angular relationships between neurons within the same layer, contributing to enhanced stability.

\section{Orthogonal Finetuning}

\subsection{Why Does Orthogonal Transformation Make Sense?}

We commence by articulating the motivations for utilizing orthogonal transformations when finetuning text-to-image diffusion models. Our analysis is divided into two fundamental aspects: (1) the rationale for adjusting the angular direction of neurons (\emph{i.e.}, their orientation), and (2) the reasons supporting the use of orthogonal transformations as a mechanism for this angular finetuning.

Addressing the first question, we are motivated by empirical findings from \cite{liu2018decoupled,chen2020angular}, which indicate that the angular difference in features provides a robust measure of the semantic gap. SphereNet~\cite{liu2017deep} demonstrates that constraining all network activations to lie on a unit hypersphere results in comparable model capacity and often enhanced generalization, suggesting that the informational content of data is well preserved in the orientation of the neurons. To further illustrate the relevance of neuron orientation, we present a simple experiment in Figure~\ref{fig:toy_ae}. Here, a conventional convolutional autoencoder is trained on a dataset of flower images. During training, feature maps are constructed via the standard inner product, with $z$ as the convolution kernel’s output for $\bm{w}$ (the kernel) and $\bm{x}$ (input patch). At inference, we evaluate three feature mapping strategies: (a) the standard inner product, (b) reliance solely on magnitude, and (c) use of only the angular component. The findings, as depicted in Figure~\ref{fig:toy_ae}, reveal that leveraging angular information alone nearly completely reconstructs the original inputs, whereas using only magnitude is uninformative. Importantly, the cosine (angular) activation in (c) is not involved during training—training employs the inner product exclusively. These observations underscore that the orientation of the neurons predominantly encodes the semantic attributes of the inputs, indicating that semantic manipulation should preferentially target adjustments to neuron directionality.

Concerning the second question, a natural approach to finetuning neuron directions is to collectively apply rotations or reflections across all neurons within a layer, inherently invoking an orthogonal transformation. Although more granular angular modifications—where each neuron is rotated independently—offer increased flexibility, our investigations suggest that orthogonal transformations optimally balance adaptability with structural discipline. Furthermore, it is established in \cite{liu2021orthogonal} that orthogonal transformations possess sufficient expressivity for neural network training. To substantiate this, we conduct a controlled generation experiment using the ControlNet~\cite{zhang2023adding} paradigm; results are illustrated in Figure~\ref{fig:ortho}. We find that the standard implementation of OFT demonstrates robust performance and precise control by training epoch 20. In contrast, omitting the orthogonality constraint in OFT causes the model to produce implausible outputs and fail in control. These outcomes reinforce the critical role of orthogonality in the efficacy of OFT.

\subsection{General Framework}

The standard approach to finetuning typically employs gradient descent with a reduced learning rate to adjust the parameters of a model or selected layers. By using a small learning rate, this process inherently restricts how far the updated model can stray from its pretrained counterpart, effectively limiting the discrepancy between the finetuned model parameters, denoted as $\bm{M}$, and the pretrained weights, $\bm{M}^0$, through the minimization of $\thickmuskip=2mu \medmuskip=2mu \| \bm{M} - \bm{M}^0\|$. Although this latent regularization is sensible, it might still provide too much latitude when refining large-scale models. To counteract this, LoRA introduces a further constraint by enforcing the weight change $\bm{M} - \bm{M}^0$ to have a low rank: $\thickmuskip=2mu \medmuskip=2mu \text{rank}(\bm{M} - \bm{M}^0)=r'$ for some small $r'$. Distinct from LoRA, OFT enforces that neuron pairwise similarity, as quantified by hyperspherical energy, remains unchanged: $\thickmuskip=2mu \medmuskip=2mu \| \text{HE}(\bm{M})-\text{HE}(\bm{M}^0) \|=0$, where $\text{HE}(\cdot)$ represents the hyperspherical energy of a matrix.

For example, consider a dense layer parameterized by $\thickmuskip=2mu \medmuskip=2mu \bm{W}=\{\bm{w}_1,\cdots,\bm{w}_n\}\in\mathbb{R}^{d\times n}$, with each column $\bm{w}_i\in\mathbb{R}^{d}$ signifying a neuron, and $\bm{W}^0$ as its corresponding pretrained weights. The layer’s output $\thickmuskip=2mu \medmuskip=2mu \bm{z}\in\mathbb{R}^n$ is computed as $\thickmuskip=2mu \medmuskip=2mu \bm{z}=\bm{W}^\top\bm{x}$ for input $\thickmuskip=2mu \medmuskip=2mu \bm{x}\in\mathbb{R}^d$. In OFT, the objective is to reduce the difference in hyperspherical energy between the adapted and the initial model weights:

\begin{equation}\label{eq:oft_principle}
\footnotesize
\min_{\bm{W}} \left\| \text{HE}(\bm{W})-\text{HE}(\bm{W}^0)\right\|~~\Leftrightarrow~~\min_{\bm{W}} \bigg{\|} \sum_{i\neq j} \|\hat{\bm{w}}_i-\hat{\bm{w}}_j\|^{-1}-\sum_{i\neq j} \|\hat{\bm{w}}^0_i-\hat{\bm{w}}^0_j\|^{-1}\bigg{\|}
\end{equation}

Here, $\thickmuskip=2mu \medmuskip=2mu \hat{\bm{w}}_i:=\bm{w}_i/\|\bm{w}_i\|$ is the normalized $i$-th neuron, and the hyperspherical energy for a weight matrix $\bm{W}$ is given by $\thickmuskip=2mu \medmuskip=2mu \text{HE}(\bm{W}):= \sum_{i\neq j} \|\hat{\bm{w}}_i-\hat{\bm{w}}_j\|^{-1}$. Notably, Eq.~\eqref{eq:oft_principle} attains its minimum value of zero when the weights $\bm{W}$ are related to the pretrained weights $\bm{W}^0$ through an orthogonal transformation, i.e., $\thickmuskip=2mu \medmuskip=2mu \bm{W}=\bm{R}\bm{W}^0$ with $\thickmuskip=2mu \medmuskip=2mu \bm{R}\in\mathbb{R}^{d\times d}$ being an orthogonal matrix (the determinant of $\bm{R}$ distinguishes between rotation and reflection, with values $1$ and $-1$, respectively). This forms the foundation of OFT

\begin{equation}\label{eq:oft_general}
    \footnotesize
    \bm{z}=\bm{W}^\top\bm{x}=(\bm{R}\cdot\bm{W}^0)^\top\bm{x},~~~\text{s.t.}~\bm{R}^\top\bm{R}=\bm{R}\bm{R}^\top=\bm{I}
\end{equation}

Here, $\bm{W}$ represents the weight matrix updated via OFT-finetuning, and $\bm{I}$ denotes the identity matrix. The overall structure of OFT is depicted in Figure~\ref{fig:block}. Analogous to the zero-initialization process adopted in LoRA, it is essential for OFT to finetune the pretrained model directly from $\bm{W}^0$. To make this possible, the orthogonal matrix $\bm{R}$ is initially set as the identity, thereby ensuring that the finetuned model begins with the pretrained parameters. To maintain the orthogonality constraint on $\bm{R}$, one can utilize differentiable orthogonalization approaches as detailed in \cite{liu2021orthogonal,lezcano2019cheap}. We will further explore how to enforce this constraint in a manner that is computationally practical.

\subsection{Efficient Orthogonal Parameterization}\label{sect:eff_ortho}

Traditional orthogonalization algorithms, such as Gram-Schmidt, are theoretically differentiable but are computationally demanding and often not suitable for large-scale applications~\cite{liu2021orthogonal}. To enhance efficiency, we leverage Cayley parameterization for the construction of the orthogonal matrix. In this formulation, the orthogonal matrix is computed as $\thickmuskip=2mu \medmuskip=2mu \bm{R}=(\bm{I}+\bm{Q})(I-\bm{Q})^{-1}$, where $\bm{Q}$ is a skew-symmetric matrix satisfying $\thickmuskip=2mu \medmuskip=2mu \bm{Q}=-\bm{Q}^\top$. This technique achieves computational efficiency, though with the minor limitation that Cayley parameterization only generates orthogonal matrices with determinant $1$ (members of the special orthogonal group). Our experiments show, however, that this constraint has no adverse effect on empirical results. Nevertheless, parameterizing $\bm{R}$ via Cayley transform still poses inefficiency when the dimension $d$ is large. To mitigate this, we propose to express $\bm{R}$ as a block-diagonal matrix with $r$ submatrices, resulting in the structure below:

\begin{equation}
    \footnotesize
    \bm{R}=\text{diag}(\bm{R}_1,\bm{R}_2,\cdots,\bm{R}_r)=\begin{bmatrix}
    \bm{R}_1\in \text{O}(\frac{d}{r}) &  &\\
    &\ddots & \\
    & & \bm{R}_r\in \text{O}(\frac{d}{r})
    \end{bmatrix}\in \text{O}(d)
\end{equation}

Here, $\text{O}(d)$ represents the orthogonal group of $d$ dimensions, with $\thickmuskip=2mu \medmuskip=2mu \bm{R}\in\mathbb{R}^{d\times d}$ and $\thickmuskip=2mu \medmuskip=2mu \bm{R}_i\in\mathbb{R}^{d/r\times d/r}$, for all $i$, as orthogonal matrices. In the special case where $\thickmuskip=2mu \medmuskip=2mu r=1$, the block-diagonal orthogonal matrix reduces to an unconstrained, full orthogonal matrix. The parameter count for a $d\times d$ orthogonal matrix is $\thickmuskip=2mu \medmuskip=2mu d(d-1)/2$, yielding $\mathcal{O}(d^2)$ complexity. If the matrix is block-diagonal with $r$ blocks, each orthogonal block, the total number of parameters becomes $\thickmuskip=2mu \medmuskip=2mu d(d/r-1)/2$, which corresponds to $\thickmuskip=2mu \medmuskip=2mu \mathcal{O}(d^2/r)$ complexity. One may reduce the parameter count even further by making all blocks identical, i.e., letting $\thickmuskip=2mu \medmuskip=2mu \bm{R}_i = \bm{R}_j$ for all $i \neq j$, thereby decreasing parameter complexity to $\mathcal{O}(d^2/r^2)$. All these variants continue to produce a matrix $\bm{R}$ that is strictly orthogonal, and thus maintain the hyperspherical energy without compromise.

Next, we analyze how OFT measures up against LoRA with respect to the number of parameters. For LoRA with rank parameter $r'$, the total number of trainable parameters is given by $\thickmuskip=2mu \medmuskip=2mu r'(d+n)$. Setting $r$ and $r'$ proportional to the neuron dimension $d$ (for example, $\thickmuskip=2mu \medmuskip=2mu r = r' = \alpha d$ where $\thickmuskip=2mu \medmuskip=2mu 0 < \alpha \leq 1$), the resulting parameter count for LoRA scales as $\mathcal{O}(d^2 + dn)$, whereas for OFT, it scales as $\mathcal{O}(d)$. To offer a direct comparison, consider a weight matrix of dimension $\thickmuskip=2mu \medmuskip=2mu 128\times128$: LoRA has $2,048$ trainable parameters when $\thickmuskip=2mu \medmuskip=2mu r'=8$, while OFT contains $960$ trainable parameters at $\thickmuskip=2mu \medmuskip=2mu r=8$, assuming no sharing among block matrices.

\subsection{Constrained Orthogonal Finetuning}

An additional restriction on the flexibility of standard OFT can be introduced by forcing the finetuned model to remain close to the original pretrained parameters. Concretely, COFT utilizes the forward propagation:

\begin{equation}\label{eq:coft_general}
    \footnotesize
    \bm{z} = \bm{W}^\top\bm{x} = (\bm{R}\cdot\bm{W}^0)^\top\bm{x},~~~\text{s.t.}~\bm{R}^\top\bm{R}=\bm{R}\bm{R}^\top=\bm{I},~\norm{\bm{R}-\bm{I}}\leq \epsilon
\end{equation}

where there are two constraints: the orthogonality condition and a deviation bound of $\epsilon$ from the identity matrix. The Cayley transform, as reviewed in Section~\ref{sect:eff_ortho}, allows enforcement of the orthogonality property. However, ensuring the additional $\epsilon$-level deviation constraint for Cayley-based orthogonal matrices poses practical difficulties. To investigate the impact of the Cayley transform, consider the Neumann series approximation: for $\thickmuskip=2mu \medmuskip=2mu \bm{R}=(\bm{I}+\bm{Q})(\bm{I}-\bm{Q})^{-1}$, this yields $\thickmuskip=2mu \medm

series converges in the operator norm). This allows us to place the constraint $\thickmuskip=2mu \medmuskip=2mu\|\bm{R}-\bm{I}\|\leq\epsilon$ directly within the Cayley transform. Accordingly, this becomes the alternative constraint $\thickmuskip=2mu \medmuskip=2mu \|\bm{Q}-\bm{0}\|\leq\epsilon'$ where $\bm{0}$ represents the zero matrix and $\epsilon'$ is a distinct tuning parameter (not equal to $\epsilon$). Enforcing this new condition on $\bm{Q}$ is straightforward by applying projected gradient descent. To initialize the orthogonal matrix $\bm{R}$ as the identity, $\bm{Q}$ is simply set to the zero matrix at initialization. Conceptually, COFT can be interpreted as integrating two concrete constraints: minimal change in hyperspherical energy, and explicit restriction on deviation from the original pretrained model. The second restriction is often only present implicitly in most finetuning approaches, while COFT elevates it to an explicit requirement. While OFT delivers strong performance, our findings indicate that COFT offers improved stability during finetuning owing to the clear imposition of the deviation limit. Figure~\ref{fig:eps_coft} illustrates the impact of varying $\epsilon$ on COFT performance. This figure shows that the value of $\epsilon$ dictates the adaptability of finetuning: a higher $\epsilon$ makes COFT's results closely match those of OFT, whereas reducing $\epsilon$ causes COFT to approximate the behaviors of the original, unmodified text-to-image diffusion model.

\subsection{Re-scaled Orthogonal Finetuning}

We introduce a straightforward modification to OFT by making a neuron-wise magnitude scaling parameter trainable. This idea arises because altering the scaling of neurons has no effect on hyperspherical energy due to normalization of vector norms. Concretely, our forward operation is given by $\thickmuskip=2mu \medmuskip=2mu\bm{z}=(\bm{R}\bm{W}^0\bm{D})^\top\bm{x}$\footnote{\textbf{Errata}: The NeurIPS camera ready manuscript incorrectly lists the re-scaled OFT forward step as $\thickmuskip=2mu \medmuskip=2mu\bm{z}=(\bm{D}\bm{R}\bm{W}^0)^\top\bm{x}$. The provided implementation is accurate; thus, results in Appendix~\ref{app:rescaled} remain valid.} where $\thickmuskip=2mu \medmuskip=2mu \bm{D}=\text{diag}(s_1,\cdots,s_n)\in\mathbb{R}^{n\times n}$ is a diagonal matrix of trainable parameters with each diagonal entry $s_1,\cdots,s_n$ constrained to be positive. Unlike the basic OFT as described in Eq.~\eqref{eq:oft_general}, where only $\bm{R}$ is optimized, in re-scaled OFT both $\bm{R}$ and the scaling matrix $\bm{D}$ are updated during training. This approach provides extra flexibility to OFT, incurring only a minimal increase in parameter count. For experimental clarity and to isolate the benefits of orthogonal transformation, we primarily employ the unscaled OFT, yet our evaluations suggest that re-scaled OFT generally yields superior results (details in Appendix~\ref{app:rescaled}).

\section{Intriguing Insights and Discussions}

\textbf{OFT’s flexibility across architectures.} The OFT method can, in principle, be applied to any neural architecture. When dealing with Transformers, LoRA is commonly utilized for adapting attention weights~\cite{hulora2022}. To ensure direct comparability with LoRA, our studies focus OFT solely on the attention weights. Besides usage in fully connected layers, OFT is also naturally amenable to adaptation of convolutional layers, especially since the block-diagonal form of $\bm{R}$ acquires appealing interpretations in this context—something LoRA lacks. By matching the number of blocks in $\bm{R}$ to the number of input channels, each block modulates a single channel independently, akin to training depth-wise convolutional filters~\cite{chollet2017xception}. Alternatively, if all blocks are tied across $\bm{R}$, the same orthogonal transformation acts across every channel, providing a global transformation. We provide an initial exploration of using OFT for convolutional layers in Appendix~\ref{app:convolution}.

\textbf{Connection to LoRA}. LoRA employs a low-rank additive matrix to ensure that updates do not drastically alter the pretrained weight matrix, preserving much of the original information. Alternatively, OFT constrains the transformation applied to the pretrained weights to be orthogonal (thus maintaining full rank), thereby safeguarding the information learned during pretraining. The forward computation in OFT can be equivalently expressed as $\thickmuskip=2mu \medmuskip=2mu \bm{z}=(\bm{R}\bm{W}^0)^\top\bm{x}=(\bm{W}^0+(\bm{R}-\bm{I})\bm{W}^0)^\top\bm{x}$, where the term $\thickmuskip=2mu \medmuskip=2mu (\bm{R}-\bm{I})\bm{W}^0$ plays a similar functional role to the low-rank adjustment made in LoRA. If $\bm{R}-\bm{I}$ is low-rank and $\bm{W}^0$ is full-rank, as is typically the case, OFT essentially introduces a low-rank modification to the weights. Both approaches have a hyperparameter controlling the parameter efficiency: LoRA utilizes the rank $r'$, while OFT employs the block-diagonal size $r$ to restrict the number of parameters updated during training. Fundamentally, LoRA and OFT exemplify two distinct strategies for achieving parameter efficiency—LoRA exploits low-rank decompositions, whereas OFT leverages sparsity arising from block-diagonal orthogonal matrices.

\textbf{Why OFT converges faster}. There is empirical evidence, as visualized in Figure~\ref{fig:toy_ae}, that changing the directionality of neurons is the most direct way to modify their semantic capacity—precisely the capability that OFT facilitates. Furthermore, OFT can be interpreted as updating weights on a smooth hyperspherical manifold, which offers a more favorable optimization landscape. Such a perspective is also supported by experimental results reported in \cite{liu2021orthogonal}.

\textbf{Why not minimize hyperspherical energy}. Distinct from the methodology in \cite{liu2021orthogonal}, we do not target minimization of hyperspherical energy. In classification tasks, it is important for neurons to avoid redundancy. The configuration corresponding to minimal hyperspherical energy enforces a uniform distribution of neurons on the hypersphere, which is not an appropriate objective during finetuning. Such uniformity risks erasing useful pretraining structure rather than adapting it.

\textbf{Balancing flexibility and regularity during finetuning}. Our findings indicate a fundamental tension between flexibility and regularity in model adaptation. While standard finetuning provides maximal flexibility, it often results in poor stability and a heightened risk of model collapse. In contrast, OFT, though conceptually straightforward, achieves a favorable compromise by maintaining the pairwise angular relationships between neurons. Moreover, employing a block-diagonal structure can be interpreted as an even stricter form of regularization imposed on the orthogonal transformation.

\textbf{Inference efficiency maintained}. In contrast to ControlNet, the OFT approach incurs zero inference-time computational cost for the adapted model. During inference, the learned orthogonal matrix $\bm{R}$ is simply merged into the pretrained parameter matrix $\bm{W}^0$, producing the effective weight matrix $\thickmuskip=2mu \medmuskip=2mu \bm{W}=\bm{R}\bm{W}^0$. Consequently, the inference runtime remains identical to that of the original pretrained architecture.

\section{Experiments and Results}

\textbf{Experimental protocol}. Our evaluation employs Stable Diffusion v1.5~\cite{rombach2022high} as the foundational text-to-image model. For objective assessment, sample images are randomly selected for each comparative approach. In the context of subject-driven image generation, we adhere to the procedures outlined in DreamBooth~\cite{ruiz2023dreambooth}. For experiments involving controllable generation, our methodology follows ControlNet~\cite{zhang2023adding} and T2I-Adapter~\cite{mou2023t2i}. To ensure methodological parity with LoRA, we restrict the application of OFT or COFT to identical layers as those employed by LoRA. Additional implementation specifics and experimental results can be found in Appendix~\ref{app:settings}.

\subsection{Subject-driven Generation}

\textbf{Experimental setup}. Baselines used include DreamBooth~\cite{ruiz2023dreambooth} and LoRA~\cite{hulora2022}. All compared methods optimize using the loss function specified in DreamBooth. For both DreamBooth and LoRA, we broadly follow their published protocols and utilize recommended hyperparameter configurations. Supplemental experimental outcomes and procedural details appear in Appendix~\ref{app:settings},~\ref{app:coft_oft},~\ref{app:more_qualitative},~\ref{app:failure}.

\textbf{Finetuning stability and convergence}. To begin, we assess the stability of finetuning as well as the convergence rate for DreamBooth, LoRA, OFT, and COFT. The corresponding outcomes are depicted in Figure~\ref{fig:teaser} and Figure~\ref{fig:db_convergence}. It is evident that both COFT and OFT can stably finetune the diffusion model throughout the process. At the 400-iteration mark, DreamBooth and OFT variants already demonstrate effective control, whereas LoRA fails to maintain the subject's identity. By 2000 iterations, DreamBooth exhibits mode collapse in its outputs, and LoRA is unable to render a yellow shirt—producing yellow fur instead. In comparison, OFT and COFT retain their capacity for stable and coherent manipulation of the generated images. These observations confirm that the OFT framework converges rapidly without sacrificing robustness in subject-driven generative tasks. Notably, robustness to finetuning iteration count is critical, as it minimizes the need for per-subject hyperparameter adjustment. With both OFT and COFT, a larger number of iterations may be chosen without demanding intensive tuning. For COFT, selecting an appropriate $\epsilon$ leads to both the learning rate and iteration count requiring little tuning effort.

\textbf{Quantitative comparison}. In alignment with the methodology of \cite{ruiz2023dreambooth}, we perform quantitative evaluations to measure subject fidelity (using DINO~\cite{caron2021emerging} and CLIP-I~\cite{radford2021learning}), alignment with text prompts (via CLIP-T~\cite{radford2021learning}), and generative diversity (assessed by LPIPS~\cite{zhang2018unreasonable}). CLIP-I is calculated as the mean pairwise cosine similarity of CLIP embeddings between real and generated images. DINO employs the same procedure as CLIP-I but substitutes CLIP with ViT S/16 DINO embeddings. CLIP-T measures the mean cosine similarity between the CLIP embeddings of generated images and their corresponding textual prompts. Diversity is evaluated through the average LPIPS cosine similarity among generated images for a specific subject and prompt. According to Table~\ref{table:dreambooth_final}, both COFT and OFT significantly surpass DreamBooth and LoRA on the DINO and CLIP-I measures, while maintaining either superior or at least comparable outcomes on text prompt fidelity and diversity metrics. To control for stochasticity, we perform each method’s finetuning procedure 30 times on the same pretrained model using distinct random seeds. These experiments demonstrate that the OFT framework excels in both convergence properties and stability, and also achieves consistently stronger overall performance.

\textbf{Qualitative comparison}. To gain a clearer, more tangible sense of the advantages provided by OFT, we present a selection of randomly chosen samples illustrating subject-driven generation in Figure~\ref{fig:dreambooth_final}. In order to ensure a level playing field, each method generates outputs using the same finetuned model, without separate prompt tuning for individual outcomes. For every technique, we utilize the checkpoint that records the highest validation CLIP score. As shown in Figure~\ref{fig:dreambooth_final}, both OFT and COFT consistently maintain semantic fidelity to the original subject, whereas LoRA frequently struggles to retain subject identity (\emph{e.g.}, the LoRA variant completely fails to preserve the subject in the bowl instance). Furthermore, OFT and COFT provide superior control over image generation via textual instructions. In contrast, DreamBooth, while capable of subject identity preservation, is less reliable at following the specified prompt (\emph{e.g.}, in the first row depicting the bowl). These visual assessments highlight OFT’s strength in balancing precise control and robust subject preservation. Additionally, OFT demonstrates robustness to the number of training iterations, maintaining stable performance even as iterations increase—a property not shared by either DreamBooth or LoRA. Expanded sets of qualitative examples are provided in Appendix~\ref{app:more_qualitative}. We also include results from a human evaluation in Appendix~\ref{app:human}, which further support the superior performance of OFT.

\subsection{Controllable Generation}

\textbf{Settings}. As baselines, we incorporate ControlNet~\cite{zhang2023adding}, T2I-Adapter~\cite{mou2023t2i}, and LoRA~\cite{hulora2022}. Three notably demanding controllable generation tasks are evaluated in the main paper: transforming Canny edges to images~(C2I) using the COCO dataset~\cite{lin2014microsoft}, converting segmentation maps to images~(S2I) with the ADE20K dataset~\cite{zhou2017scene}, and generating faces from landmarks~(L2F) on CelebA-HQ~\cite{karras2018progressive,xia2021tedigan}. Each method fine-tunes Stable Diffusion~(SD) v1.5 for 20 epochs on these datasets. Additional experimental results are detailed in Appendix~\ref{app:more_qualitative},\ref{app:more_control_task},\ref{app:failure}.

\textbf{Convergence}. To assess convergence rates, we compare ControlNet, T2I-Adapter, LoRA, and COFT on the L2F task, providing both numerical and visual analyses. For our quantitative metric, we calculate the average $\ell_2$ norm between the intended control face landmarks and those generated by the models. Figure~\ref{fig:face_convergence} presents how the landmark error decreases as models are finetuned over different epochs. It is evident from these results that COFT and OFT reach low error values substantially more quickly than alternatives; for example, LoRA needs 20 epochs to attain the same performance level that OFT reaches at epoch 8. It is noteworthy that OFT and COFT involve a comparable number of trainable parameters as LoRA (and substantially fewer than ControlNet), yet their convergence is much more rapid than these prior methods. Additionally, the expedited convergence of OFT is further supported by Figure~\ref{fig:teaser}, whose rightmost example highlights OFT’s superior data efficiency relative to ControlNet and LoRA: OFT is capable of converging effectively using just 5\% of the ADE20K dataset. In qualitative comparisons, we primarily examine OFT, COFT, and ControlNet, since ControlNet most closely matches OFT’s landmark accuracy. Figure~\ref{fig:face_convergence_qual} demonstrates that both OFT and COFT not only converge stably but also progressively adjust the generated facial pose to correspond with the target landmarks. Contrasting with OFT’s consistent and smooth trajectory, ControlNet exhibits an abrupt onset of controllability after epoch 8, mirroring the sudden convergence behavior reported in \cite{zhang2023adding}. The overall stability of OFT’s convergence greatly simplifies practical application, as its training progression is markedly more consistent and foreseeable—making it more straightforward to identify optimal hyperparameters.

\textbf{Quantitative comparison}. To assess the effectiveness of controllable generation, we propose a control consistency metric. This metric is based on extracting the control signal from the output image and directly comparing it to the original input control. For the C2I setting, we report IoU and F1 scores, whereas for S2I, we provide mean IoU, as well as mean and overall accuracy. The L2F evaluation uses the mean $\ell_2$ distance between control landmarks and those predicted by the model. Comprehensive descriptions of these consistency measurements are provided in Appendix~\ref{app:settings}. Each baseline is optimized using its best possible hyperparameter configuration. As detailed in Table~\ref{table:control_gen}, both OFT and COFT consistently demonstrate superior and more precise control capabilities in comparison to alternative methods. In contrast, approaches that rely on adapters (\emph{e.g.}, T2I-Adapter and ControlNet) show slower convergence and ultimately produce inferior results. While LoRA surpasses ControlNet on the S2I benchmark, it underperforms on C2I and L2F, and overall, its potential appears limited, despite extensive tuning. Notably, our experiments indicate that the OFT framework continues to improve as the trainable parameter count increases, without clear signs of plateauing. It is important to highlight that our rigorous quantitative approach to evaluating controllable generation represents a novel aspect of our work, as it enables precise assessment of control performance for finetuned models, bounded only by the accuracy of the underlying segmentation or detection models employed.

\textbf{Qualitative comparison}. We further conduct a qualitative assessment between OFT, COFT, and leading methods such as ControlNet, T2I-Adapter, and LoRA. As illustrated by the randomly sampled outputs in Figure~\ref{fig:control_final}, OFT and COFT consistently generate images with superior realism and fidelity while maintaining precise controllability. Specifically, for the S2I task, LoRA fails to produce images adhering to the provided segmentation maps; in contrast, ControlNet, OFT, and COFT demonstrate robust control over image synthesis. Relative to ControlNet, both OFT and COFT generate images with sharper details and more coherent geometric structures, all while utilizing significantly fewer parameters. Regarding the C2I task, OFT and COFT successfully synthesize semantically relevant visuals from coarse Canny edge inputs, whereas T2I-Adapter and LoRA show markedly inferior performance. For the L2F task, our approach yields the most reliable pose consistency in generated facial images, even under challenging orientation conditions. Across all three controlled generation tasks, OFT and COFT provide superior qualitative outcomes compared to state-of-the-art alternatives, highlighting the advantages of our OFT-based framework. Additional qualitative results for each control scenario are available in Appendix~\ref{app:control_exp}. Furthermore, Appendix~\ref{app:more_control_task} contains demonstrations of OFT's capabilities across further control domains, such as dense pose to body, sketch to image, and depth to image translation.

\section{Concluding Remarks and Open Problems}

Drawing inspiration from the critical role that angular relationships among neurons play in defining visual semantics, we introduce a straightforward yet robust finetuning strategy—orthogonal finetuning—for managing text-to-image diffusion models. Our approach specifically addresses two domains: subject-driven generation and controllable generation. In contrast to current alternatives, OFT achieves enhanced controllability and improved stability during finetuning, while requiring a reduced set of trainable parameters. Additionally, OFT does not incur any extra computational burden during inference, rendering it a practical and efficient solution for deployment.

OFT also raises several notable open questions. First, OFT enforces orthogonality through Cayley parametrization, necessitating a matrix inversion operation that can impede scalability. We partially resolve this concern by implementing a block diagonal parametrization; however, developing faster, differentiable techniques for matrix inversion remains unresolved. Second, the compositional nature of OFT suggests promising avenues for future work: the orthogonal matrices generated from multiple OFT finetuning instances can be composed via multiplication to yield another orthogonal matrix. Investigating whether this compositional property allows for the retention of knowledge across diverse downstream tasks is an open question. Finally, the model’s parameter efficiency is chiefly governed by the use of block diagonal structures, which can introduce systematic biases and restrict overall flexibility. Identifying more effective, unbiased strategies for increasing parameter efficiency represents another pressing open challenge.

\end{document}